\documentclass[11pt,reqno]{amsart}

%==============================================================
%  An algebraic--variational framework for the affine plank problem
%  Working manuscript. Genre: framework + research programme.
%
%  ACADEMIC-CORRECTNESS POLICY (binding):
%  Every environment named "Theorem", "Proposition", "Lemma", or
%  "Corollary" is EITHER (a) a result from the cited literature,
%  restated with attribution, OR (b) an elementary statement proved
%  in full here. Every not-yet-established statement is placed in a
%  "Conjecture", "Problem", or "Programme" environment and is never
%  used as if proved. See Section 1.3 (Status of results).
%==============================================================

\usepackage{amsmath,amssymb,amsthm}
\usepackage[margin=1.1in]{geometry}
\usepackage{mathtools}
\usepackage[colorlinks=true,linkcolor=blue,citecolor=blue]{hyperref}
\usepackage[protrusion=true,expansion=false]{microtype}

\theoremstyle{plain}
\newtheorem{theorem}{Theorem}[section]
\newtheorem{proposition}[theorem]{Proposition}
\newtheorem{lemma}[theorem]{Lemma}
\newtheorem{corollary}[theorem]{Corollary}

\theoremstyle{definition}
\newtheorem{conjecture}[theorem]{Conjecture}
\newtheorem{problem}[theorem]{Problem}
\newtheorem{definition}[theorem]{Definition}
\newtheorem{remark}[theorem]{Remark}
\newtheorem{example}[theorem]{Example}

\newcommand{\R}{\mathbb{R}}
\newcommand{\C}{\mathbb{C}}
\newcommand{\N}{\mathbb{N}}
\newcommand{\inner}[2]{\langle #1, #2\rangle}
\newcommand{\norm}[1]{\lVert #1\rVert}
\newcommand{\abs}[1]{\lvert #1\rvert}
\newcommand{\relw}{\operatorname{rw}}
\newcommand{\wid}{\operatorname{w}}
\newcommand{\sgn}{\operatorname{sgn}}
\newcommand{\diag}{\operatorname{diag}}
\newcommand{\Jac}{\operatorname{J}}
\DeclareMathOperator{\spann}{span}

\begin{document}

\title[An algebraic--variational framework for the affine plank problem]
{An algebraic--variational framework for the affine plank problem}

\author{}
\address{}
\email{}

\date{Working draft, \today}

\subjclass[2020]{Primary 52A40, 52C17; Secondary 46B07, 14M25, 90C26.}
\keywords{plank problem, relative width, affine plank conjecture, polarization
constants, Euler--Jacobi vanishing theorem, Bernstein--Kushnirenko bound,
critical points.}

\begin{abstract}
Bang's affine (relative--width) plank conjecture asserts that if a convex body
$K\subset\R^d$ is covered by finitely many planks, then the sum of the widths of
the planks, each measured relative to the width of $K$ in the corresponding
direction, is at least $1$. The conjecture is sharp, is known for centrally
symmetric bodies by a theorem of Ball, and is otherwise open---open even in the
plane and even for three planks. The best current lower bound for the total
relative width of a covering of a general body is $2/(1+\sqrt d)$, due to Bakaev
and Yehudayoff. We isolate a methodological gap: the critical--point and
Euler--Jacobi vanishing technique that Mart\'inez and Ortega--Moreno recently
used to settle the real linear polarization problem has not been brought to bear
on planks. This paper sets up a precise framework for doing so. We recall the
engine in the polarization setting, record the elementary dictionary between
relative widths and gauge functionals, observe that in the centrally symmetric
(Hilbertian) case the relative--width statement is exactly of polarization type
and is therefore already accessible to the engine, and formulate the three
analytic obstructions---reality of the critical locus, a Bernstein--Kushnirenko
count of that locus, and a toric Euler--Jacobi vanishing identity---as explicit
conjectures whose resolution would transport the method to the non--symmetric
case. We are careful throughout to separate what is proved from what is proposed.
\end{abstract}

\maketitle

\section{Introduction}\label{sec:intro}

\subsection{The plank problem and its affine refinement}
A \emph{plank} of width $w\ge 0$ in $\R^d$ is the closed region between two
parallel hyperplanes at Euclidean distance $w$; equivalently, for a unit vector
$u\in S^{d-1}$ and reals $a\le b$ with $b-a=w$,
\[
P=\{x\in\R^d:\ a\le \inner{x}{u}\le b\}.
\]
For a convex body $C\subset\R^d$ (a compact convex set with nonempty interior),
the \emph{width of $C$ in the direction $u$} is
\[
\wid(C;u)=\max_{x\in C}\inner{x}{u}-\min_{x\in C}\inner{x}{u},
\]
and the \emph{minimal width} of $C$ is $\wid(C)=\min_{u\in S^{d-1}}\wid(C;u)$.

In 1932 Tarski asked whether a convex body covered by planks must have the sum of
the plank widths at least $\wid(C)$; he settled only the disk
\cite{Tarski1932}. Bang proved Tarski's conjecture in general
\cite{Bang1951}: if $C\subseteq P_1\cup\dots\cup P_n$ then
$\sum_i \wid(P_i)\ge \wid(C)$. At the end of \cite{Bang1951} Bang proposed a
stronger, affinely invariant version, introducing the term \emph{relative width}.

\begin{definition}\label{def:relw}
Let $P$ be a plank with unit normal $u$ and width $\wid(P)$. The \emph{relative
width} of $P$ with respect to a convex body $C$ is
\[
\relw_C(P)\ :=\ \frac{\wid(P)}{\wid(C;u)}\,.
\]
\end{definition}

\begin{conjecture}[Bang's affine plank conjecture {\cite{Bang1951}}]
\label{conj:bang}
If a convex body $C\subset\R^d$ ($d\ge2$) is covered by planks
$P_1,\dots,P_n$, then
\[
\sum_{i=1}^{n}\relw_C(P_i)\ \ge\ 1 .
\]
\end{conjecture}

The quantity $\sum_i\relw_C(P_i)$ is invariant under invertible affine maps of
$\R^d$, which is the sense in which Conjecture~\ref{conj:bang} refines Tarski's
absolute--width problem. The bound is sharp: $d$ parallel slabs of equal relative
width $1/d$ stacked across $C$ already attain equality in many configurations.

\subsection{What is known}\label{subsec:known}
We summarize the state of the art; precise quotations and a verification log are
collected in the companion notes accompanying this manuscript.

\begin{itemize}
\item \textbf{Symmetric case (solved).} Ball \cite{Ball1991} proved
Conjecture~\ref{conj:bang} for centrally symmetric bodies. In functional--analytic
form: if the unit ball of a Banach space is covered by a countable family of
planks, the sum of their half--widths is at least $1$; dually, for a symmetric
convex body $C$ and $n$ hyperplanes there is a translate of a copy of $C$ scaled
by at least $1/(n+1)$ lying inside $C$ and avoiding all the hyperplanes. The proof
extends Bang's Lemma and turns on the optimization of a suitable quadratic form
over the sign vectors $\{-1,1\}^n$ of a symmetric matrix.

\item \textbf{General case (open).} For arbitrary convex bodies
Conjecture~\ref{conj:bang} is open. It is open already in the plane, and is not
known even for three planks; Bang himself verified only the case of two planks
\cite{Bang1951,KupavskiiPach2017,Verreault2026,BakaevYehudayoff2026}. The known
lower bounds for the total relative width of a covering of a general body in
$\R^d$ have improved from $1/d$ (John's theorem together with Bang's theorem; see
\cite{BezdekLitvak2009,ChambersMouille2016}) to $2/(1+d)$
\cite{ChambersMouille2016} and, most recently, to $2/(1+\sqrt d)$
\cite{BakaevYehudayoff2026}, the current record. Chambers and Mouille
\cite{ChambersMouille2016}
also prove the unconditional bound $\sum_i\relw_C(P_i)\ge 1/\dim\pi_V(C)$, where
$V$ is the span of the plank normals and $\pi_V$ is orthogonal projection onto
$V$, and they prove the sharp bound $1$ \emph{conditionally}, under the hypothesis
that $C$ minus the union of the first $m$ planks remains convex for every $m$.
\end{itemize}

\subsection{The polarization connection and a methodological gap}
\label{subsec:polar}
The plank problem is tightly linked to lower bounds for products of linear
functionals. For a Banach space $X$, the $n$th linear polarization constant
$c_n(X)$ is the least $C$ with $\norm{f_1}\cdots\norm{f_n}\le C\,\norm{f_1\cdots
f_n}$ for all $f_1,\dots,f_n\in X^{*}$; Ball's symmetric plank theorem yields the
general real bound $c_n(X)\le n^{n}$
\cite{BenitezSarantopoulosTonge1998,ReveszSarantopoulos2004}. Over a real Hilbert
space the sharp conjecture $c_n\le n^{n/2}$ stood for decades and was recently
established. Mart\'inez and Ortega--Moreno \cite{MartinezOrtegaMoreno2026} solved
the real polarization problem by a variational analysis of the critical locus of
a product of linear forms on the sphere, combined with the classical
\emph{Euler--Jacobi vanishing theorem}. Galicer, Ortega--Moreno and Pinasco
\cite{GalicerOrtegaPinasco2026} extended the argument to a weighted ``strong''
inequality and, as a corollary, to a centred plank theorem in Hilbert space.

These developments introduce into the plank/polarization circle a genuinely
algebraic tool---counting the solutions of a polynomial critical system and
forcing a good one through a vanishing identity. To the best of our knowledge,
\emph{this tool has not been applied to the affine plank problem}. The purpose of
the present paper is to set up, precisely and honestly, a framework that would do
so, and to delineate exactly which analytic facts must be established for the
transfer to succeed.

\subsection{Status of results}\label{subsec:status}
In keeping with the affine, sharp, and still--open nature of
Conjecture~\ref{conj:bang}, we are explicit about what this manuscript does and
does not prove.
\begin{itemize}
\item Section~\ref{sec:engine} \emph{recalls}, with attribution, the
critical--point/Euler--Jacobi engine of
\cite{MartinezOrtegaMoreno2026,GalicerOrtegaPinasco2026}. Nothing there is claimed
as new.
\item Section~\ref{sec:dictionary} proves elementary, self--contained facts: the
affine invariance of the total relative width
(Proposition~\ref{prop:affineinv}), the gauge reformulation
(Proposition~\ref{prop:gauge}), and the observation that in the centrally
symmetric case the relevant statement is of polarization type and hence already
within reach of the engine (Proposition~\ref{prop:symmetric}). These are correct
and complete.
\item Section~\ref{sec:programme} states the three transfer obstructions as
\emph{conjectures and problems}. They are \emph{not} theorems, are clearly labeled
as such, and are never used as if proved.
\end{itemize}
No new lower bound for the general case, and in particular no improvement of
$2/(1+\sqrt d)$, is claimed here. The contribution of this paper is a framework
and a programme, together with the elementary results that make the programme
precise.

\section{The critical--point/Euler--Jacobi engine, recalled}\label{sec:engine}

We recall the mechanism of \cite{MartinezOrtegaMoreno2026} in the streamlined form
of \cite{GalicerOrtegaPinasco2026}. Nothing in this section is original; we
include it because the framework of Section~\ref{sec:programme} is a structural
analogy with it.

Let $v_1,\dots,v_n$ be a basis of $\R^n$ and let $k_1,\dots,k_n\in\N$,
$s=\sum_j k_j$, $\alpha_j=k_j/s$. Consider
\[
P(x)=\prod_{j=1}^{n}\inner{v_j}{x}^{k_j}, \qquad x\in S^{n-1},
\]
and write $y_j=\inner{v_j}{x}$. With $w_1,\dots,w_n$ the dual basis and
$z_j=\inner{w_j}{x}$ one has $x=\sum_j z_j v_j$ and $z=Ay$, where $A=G^{-1}$ and
$G=(\inner{v_i}{v_j})_{i,j}$ is the Gram matrix.

\begin{lemma}[Critical system; \cite{MartinezOrtegaMoreno2026,GalicerOrtegaPinasco2026}]
\label{lem:critsystem}
At a critical point of $\log\abs{P}$ on the sphere with $P\neq0$ one has
$\sum_j k_j v_j/y_j=s\,x$ and $y_j z_j=\alpha_j$. Equivalently, the critical
points are the solutions of the quadratic system $h(y)=0$, where
\[
h_j(y)=y_j (Ay)_j-\alpha_j,\qquad j=1,\dots,n .
\]
\end{lemma}

\begin{lemma}[Reality and count;
{\cite[Lem.~3--4]{GalicerOrtegaPinasco2026}},
{\cite{MartinezOrtegaMoreno2026}}]\label{lem:realcount}
Every solution $y\in\C^n$ of $h(y)=0$ is real; in each open sign orthant of
$\R^n$ there is exactly one solution; hence $h$ has exactly $2^n$ solutions, all
real and simple. At a solution,
$\Jac h(y)=\diag(y)\big(A+\diag(\alpha_j/y_j^2)\big)$, which is invertible.
\end{lemma}

\begin{theorem}[Polarization via Euler--Jacobi;
{\cite{MartinezOrtegaMoreno2026,GalicerOrtegaPinasco2026}}]\label{thm:euler}
With the notation above there is a unit vector $u\in S^{n-1}$ with
$\sum_j k_j^2/\inner{v_j}{u}^2\le s^2$, equivalently
$\sum_j \alpha_j^2/\inner{v_j}{u}^2\le 1$.
\end{theorem}

\begin{proof}[Sketch, after {\cite{GalicerOrtegaPinasco2026}}]
The system $h=0$ is a zero--dimensional complete intersection of $n$ quadrics with
$2^n=\prod_j\deg h_j$ simple solutions (Lemma~\ref{lem:realcount}), so B\'ezout's
bound is attained and the Euler--Jacobi vanishing theorem applies: for every
polynomial $g$ with $\deg g\le n-1$,
\[
\sum_{h(y)=0}\frac{g(y)}{\det \Jac h(y)}=0 .
\]
Choosing $g(y)=Q(y)\big(s^2-\sum_j k_j^2/y_j^2\big)$ with $Q(y)=\prod_j y_j$, and
using $\det \Jac h(y)=Q(y)\det\!\big(A+\diag(\alpha_j/y_j^2)\big)$ with the second
factor positive at every solution, the identity becomes a sum of the quantities
$s^2-\sum_j k_j^2/y_j^2$ weighted by positive numbers. Hence some solution has
$\sum_j k_j^2/y_j^2\le s^2$; its normalization is the desired $u$. Rational and
real weights follow by approximation; see
\cite{GalicerOrtegaPinasco2026}.
\end{proof}

\begin{remark}\label{rem:threepillars}
Theorem~\ref{thm:euler} rests on three pillars: \emph{(P1) reality} of the
critical locus, \emph{(P2) an exact count} of that locus matching the product of
the degrees, and \emph{(P3) a vanishing identity} forcing a witness with the
desired sign. The metric data enter only through the positive--definite matrix
$A=G^{-1}$. The framework of Section~\ref{sec:programme} asks what survives when
the round sphere is replaced by the boundary of a general convex body and $A$ by
the second--order data of that body.
\end{remark}

\section{Elementary reformulations}\label{sec:dictionary}

This section is self--contained and its results are proved in full.

\subsection{Affine invariance}
\begin{proposition}\label{prop:affineinv}
Let $T:\R^d\to\R^d$ be an invertible affine map. If planks $P_1,\dots,P_n$ cover a
convex body $C$, then $T(P_1),\dots,T(P_n)$ are planks covering $T(C)$ and
\[
\sum_{i=1}^n\relw_{T(C)}(T(P_i))=\sum_{i=1}^n\relw_{C}(P_i).
\]
\end{proposition}

\begin{proof}
Write $T(x)=Lx+b$ with $L$ invertible. The image of a plank with normal $u$ and
width $w$ is the region between the images of its two boundary hyperplanes; these
remain parallel, so $T(P_i)$ is a plank, with some normal $u_i'$ and width
$w_i'$. For any direction, affine images of the two supporting hyperplanes of
$C$ that realize $\wid(C;u)$ are supporting hyperplanes of $T(C)$ parallel to the
images of the bounding hyperplanes of $P_i$. Both $\wid(T(P_i))$ and
$\wid(T(C);u_i')$ are the distances between the corresponding pairs of parallel
hyperplanes, measured along the same transversal; their ratio is therefore
unchanged because it is computed within a single one--dimensional quotient
$\R^d/(\text{common hyperplane direction})$ on which $T$ acts by a single affine
isomorphism. Hence $\relw_{T(C)}(T(P_i))=\relw_{C}(P_i)$ for each $i$, and the
covering $T(C)\subseteq\bigcup_i T(P_i)$ is immediate from
$C\subseteq\bigcup_i P_i$.
\end{proof}

Proposition~\ref{prop:affineinv} is the reason Conjecture~\ref{conj:bang} is
called \emph{affine}: one may normalize $C$ by any affine map (for example into
John position) without changing the quantity under study.

\subsection{Gauge reformulation}
For a convex body $C$ with $0$ in its interior, let
$\norm{x}_C=\inf\{t>0: x\in tC\}$ be its gauge (Minkowski functional) and
$h_C(u)=\max_{x\in C}\inner{x}{u}$ its support function. For a plank $P$ with unit
normal $u$ given by $a\le\inner{\cdot}{u}\le b$, one has
$\wid(P)=b-a$ and $\wid(C;u)=h_C(u)+h_C(-u)$.

\begin{proposition}\label{prop:gauge}
With $C$, $P_i$ (normals $u_i$, levels $a_i\le b_i$) as above,
\[
\relw_C(P_i)=\frac{b_i-a_i}{h_C(u_i)+h_C(-u_i)} .
\]
In particular, if $C$ is centrally symmetric about $0$ then
$h_C(u)+h_C(-u)=2h_C(u)=2\norm{u}_{C^{\circ}}$, where $C^\circ$ is the polar body,
and $\relw_C(P_i)=(b_i-a_i)/\big(2\norm{u_i}_{C^\circ}\big)$.
\end{proposition}

\begin{proof}
The first identity is the definition together with
$\wid(C;u)=h_C(u)+h_C(-u)$. If $C=-C$ then $h_C(-u)=h_C(u)$, and
$h_C=\norm{\cdot}_{C^\circ}$ by the definition of the polar body $C^\circ=\{u:
\inner{x}{u}\le1\ \forall x\in C\}$.
\end{proof}

\subsection{The symmetric case is of polarization type}
The next proposition records that, for symmetric bodies, the relative--width
statement is exactly the kind of inequality the engine of
Section~\ref{sec:engine} produces. It is a reading of known results, included to
mark the boundary of what is already accessible.

\begin{proposition}\label{prop:symmetric}
Let $H$ be a real Hilbert space (so $C=B_H$ is the symmetric Euclidean ball).
A family of planks with unit normals $u_i$ and half--widths $t_i\ge 0$ fails to
cover $B_H$ as soon as $\sum_i t_i<1$. Equivalently, in the dual formulation of
\cite{Ball1991,GalicerOrtegaPinasco2026}: given unit vectors $v_1,\dots,v_n\in H$
and weights $p_i>0$ with $\sum_i p_i=1$, there is a unit vector $u$ with
$\abs{\inner{v_i}{u}}\ge p_i$ for all $i$; in particular the planks
$\{x:\abs{\inner{v_i}{x}}<p_i\}$ do not cover the unit sphere.
\end{proposition}

\begin{proof}
This is the centred Hilbert plank statement. The displayed dual form is
\cite[Cor.~6]{GalicerOrtegaPinasco2026}, itself a consequence of
Theorem~\ref{thm:euler} via the weighted arithmetic--geometric mean inequality;
the covering statement is the contrapositive. The symmetric body case of
Conjecture~\ref{conj:bang} in full generality (countable coverings, sharp constant)
is Ball's theorem \cite{Ball1991}.
\end{proof}

\begin{remark}
Proposition~\ref{prop:symmetric} delineates the frontier precisely: for symmetric
$C$ the engine already applies (through the Gram matrix $A=G^{-1}$ of the
normals, exactly as in Section~\ref{sec:engine}). The open content of
Conjecture~\ref{conj:bang} is entirely in the \emph{non--symmetric} case, where
$h_C(u)+h_C(-u)$ is genuinely asymmetric and no inner--product structure is given.
\end{remark}

\section{A programme for the non--symmetric case}\label{sec:programme}

We now state, as conjectures and problems, the analytic facts whose proof would
transport the engine of Section~\ref{sec:engine} to general bodies. We stress that
this section contains \emph{no theorems}: the statements below are open and are
offered as a research programme.

\subsection{The proposed critical system}
Fix a smooth, strictly convex body $C$ with $0$ in its interior and a covering by
planks with normals $u_1,\dots,u_n$. Guided by Section~\ref{sec:engine}, one seeks
an objective $\Phi$ on $\partial C$ (or on a suitable gauge sphere) whose critical
points encode the total relative width, with the role of the Gram matrix played by
the second fundamental form / support--function Hessian of $C$. Making this precise
is the first task.

\begin{problem}\label{prob:objective}
Construct an objective functional $\Phi$, polynomial or rational in suitable
coordinates adapted to $\partial C$, whose stationarity conditions on
$\partial C$ form a zero--dimensional polynomial system $H_C(y)=0$ encoding the
quantities $\relw_C(P_i)$, and which reduces to the system $h$ of
Lemma~\ref{lem:critsystem} when $C$ is an ellipsoid.
\end{problem}

\subsection{The three obstructions}
Assume Problem~\ref{prob:objective} is solved, producing a system $H_C$ with
second--order data encoded by a positive--definite matrix field $A_C$ (the inverse
support--Hessian).

\begin{conjecture}[P1: reality]\label{conj:reality}
For smooth, strictly convex $C$ every complex solution of $H_C(y)=0$ is real.
\end{conjecture}

The motivation is that the reality argument of
\cite[Lem.~3]{GalicerOrtegaPinasco2026} uses only positive--definiteness of $A$
through an imaginary--part energy identity; for a smooth strictly convex body the
support--Hessian is positive definite, so the argument is expected to persist,
possibly after a limiting passage to handle bodies that are not strictly convex
(such as polytopes).

\begin{conjecture}[P2: BKK count]\label{conj:count}
The system $H_C(y)=0$ is a zero-dimensional complete intersection. Its number of
solutions equals the Bernstein--Kushnirenko (BKK) mixed volume of the Newton
polytopes of its components; all solutions are simple; and they are in bijection
with the normal-fan cells of $C$ cut out by the plank directions.
\end{conjecture}

For the round sphere the Newton polytopes are those of the quadrics $h_j$ and the
BKK number specializes to the B\'ezout number $2^n$ of Lemma~\ref{lem:realcount};
Conjecture~\ref{conj:count} asks for the analogue when the body is not an
ellipsoid, where the mixed volume is in general smaller and structured.

\begin{conjecture}[P3: toric Euler--Jacobi vanishing]\label{conj:vanishing}
There is a test functional $g_C$, of degree below the relevant BKK threshold, for
which the toric (global residue) form of the Euler--Jacobi vanishing theorem
yields
\[
\sum_{H_C(y)=0}\frac{g_C(y)}{\det \Jac H_C(y)}=0,
\]
and $g_C$ can be chosen so that this identity forces the existence of a critical
point witnessing $\sum_i\relw_C(P_i)\ge 1$.
\end{conjecture}

\medskip
\noindent\textbf{Intended architecture (conditional, not a theorem).}
\emph{If} Problem~\ref{prob:objective} is solved and
Conjectures~\ref{conj:reality}--\ref{conj:vanishing} hold for a class
$\mathcal{K}$ of bodies, \emph{then} Conjecture~\ref{conj:bang} would follow for
every $C\in\mathcal{K}$. We state this only to display the intended logical
architecture; it is contingent on three open conjectures and is deliberately
\emph{not} formulated as a theorem.
\medskip

\subsection{Test classes and falsifiers}
The programme is to be validated or refuted on the smallest non--trivial classes
before any attempt at generality.

\begin{problem}[Symmetric recovery]\label{prob:M1}
Carry out Problem~\ref{prob:objective} and
Conjectures~\ref{conj:reality}--\ref{conj:vanishing} for centrally symmetric
$C$, thereby \emph{reproving} Ball's theorem \cite{Ball1991} by the
critical--point/Euler--Jacobi method. (By Proposition~\ref{prop:symmetric} the
required inputs are already available; the content is to phrase Ball's quadratic
optimization as a special case of the engine.)
\end{problem}

\begin{problem}[Simplices]\label{prob:M2}
Establish the sharp bound $\sum_i\relw_C(P_i)\ge1$ for the simplex, the most
tractable non--symmetric body. If a reduction of Conjecture~\ref{conj:bang} to
simplices in the literature is correct, a positive answer here would settle the
conjecture in a structured case for the first time by these methods.
\end{problem}

\begin{problem}[Planar three--plank case]\label{prob:M3b}
Decide Conjecture~\ref{conj:bang} for $d=2$ and $n=3$, where the critical system
is of small size, by the present method. This case is currently open
\cite{Verreault2026,KupavskiiPach2017}.
\end{problem}

A negative outcome on Problem~\ref{prob:M1} (the engine failing to recover even
the symmetric, already--known case) would be decisive evidence against the
algebraic route and would redirect the programme toward purely variational, non
algebraic, formulations.

\section*{Acknowledgements}
(To be added.)

\begin{thebibliography}{99}

\bibitem{BakaevYehudayoff2026}
E.~Bakaev and A.~Yehudayoff,
\emph{A note on the affine plank conjecture},
preprint, arXiv:2602.20290, 2026.

\bibitem{Ball1991}
K.~M.~Ball,
\emph{The plank problem for symmetric bodies},
Invent. Math. \textbf{104} (1991), 535--543.

\bibitem{Ball2001}
K.~M.~Ball,
\emph{The complex plank problem},
Bull. London Math. Soc. \textbf{33} (2001), 433--442.

\bibitem{Bang1951}
T.~Bang,
\emph{A solution of the ``plank problem''},
Proc. Amer. Math. Soc. \textbf{2} (1951), 990--993.

\bibitem{BenitezSarantopoulosTonge1998}
C.~Ben\'itez, Y.~Sarantopoulos and A.~Tonge,
\emph{Lower bounds for norms of products of polynomials},
Math. Proc. Cambridge Philos. Soc. \textbf{124} (1998), 395--408.

\bibitem{BezdekLitvak2009}
K.~Bezdek and A.~E.~Litvak,
\emph{Covering convex bodies by cylinders and lattice points by flats},
J. Geom. Anal. \textbf{19} (2009), 233--243.

\bibitem{ChambersMouille2016}
G.~R.~Chambers and S.~Mouille,
\emph{A note on the affine--invariant plank problem},
preprint, arXiv:1604.00456, 2016.

\bibitem{GalicerOrtegaPinasco2026}
D.~Galicer, O.~Ortega--Moreno and D.~Pinasco,
\emph{Strong polarization and entropy},
preprint, arXiv:2606.02567, 2026.

\bibitem{KupavskiiPach2017}
A.~Kupavskii and J.~Pach,
\emph{From Tarski's plank problem to simultaneous approximation},
Amer. Math. Monthly \textbf{124} (2017), no.~6, 494--505.

\bibitem{MartinezOrtegaMoreno2026}
A.~D.~Mart\'inez and O.~Ortega--Moreno,
\emph{A solution to the polarization problem},
preprint, 2026.

\bibitem{ReveszSarantopoulos2004}
S.~G.~R\'ev\'esz and Y.~Sarantopoulos,
\emph{Plank problems, polarization and Chebyshev constants},
J. Korean Math. Soc. \textbf{41} (2004), no.~1, 157--174.

\bibitem{Tarski1932}
A.~Tarski,
\emph{Remarks on the degree of equivalence of polygons} (Polish),
Parametr \textbf{2} (1932), 310--314.

\bibitem{Verreault2026}
W.~Verreault,
\emph{Plank theorems and their applications: a survey},
Bull. London Math. Soc. (2026); preprint arXiv:2203.05540.

\end{thebibliography}

\end{document}
